\documentclass{article}
\usepackage{amsmath}
\usepackage{amssymb}

\title{Homework 3}
\author{Brandon Reich}
\date{30 January 2026}

\begin{document}

\maketitle

\subsection*{2.1.1}
\begin{itemize}
    \item b - Valid
    $$
        \begin{array}{|c|c|c|c|}
            \hline
            p & q & p \leftrightarrow q & p \lor q \\
            \hline
            T & T & T & T \\
            T & F & F & T \\
            F & T & F & T \\
            F & F & T & F \\
            \hline
        \end{array}
    $$
    \item g - Invalid, $p=q=F$
\end{itemize}

\subsection*{2.1.3}
\begin{itemize}
    \item c - Invalid, $p=F, q=T$ \\
    
    $p: \sqrt{2}$ is an irrational number \\
    $q: 2\sqrt{2}$ is an irrational number
    $$
        \begin{array}{c}
            \text{Form of the argument} \\
            \hline \\
            p \rightarrow q \\
            q \\
            \hline 
            \therefore p \\
        \end{array}
    $$
\end{itemize}

\subsection*{2.2.1}
\begin{itemize}
    \item a - Invalid, $\text{Sally had a side effect}=T, \text{Sally took the medication}=T$
    \item d - Valid, Modus Tollens
\end{itemize}

\subsection*{2.2.2}
\begin{itemize}
    \item f
    $$
        \begin{array}{|c|c|c|}
            \hline
            1. & p \land r & \text{Hypothesis} \\
            \hline
            2. & p & \text{Simplification, } 1 \\
            \hline
            3. & r & \text{Simplification, } 1 \\
            \hline
            4. & p \rightarrow q & \text{Hypothesis} \\
            \hline
            5. & q & \text{Modus pones, } 2,4 \\
            \hline
            6. & r \rightarrow u & \text{Hypothesis} \\
            \hline
            7. & u & \text{Modusß pones, } 3,6 \\
            \hline
            8. & q \land u & \text{Conjunction, } 5,7 \\
            \hline
        \end{array}
    $$
\end{itemize}

\subsection*{2.2.4}
\begin{itemize}
    \item a - Valid \\
    \\
    $d:$ I will drive on the freeway \\
    $f:$ I will see fire \\
    $s:$ I will take the side streets
    
    $$
    \begin{array}{c}
        \text{Form of the argument} \\
        \hline \\
        d \rightarrow f \\
        d \lor s \\
        \neg s \\
        \hline
        \therefore f
    \end{array}
    $$
    \\

    \begin{center}
    Proof
    \end{center}
    $$
    \begin{array}{|c|c|c|}
        \hline
        1. & d \lor s & \text{Hypothesis} \\
        \hline
        2. & \neg s & \text{Hypothesis} \\
        \hline
        3. & d & \text{Disjunctive syllogism } 1,2 \\
        \hline
        4. & d \rightarrow f & \text{Hypothesis} \\
        \hline
        5. & f & \text{Modus pones } 3,4 \\
        \hline
    \end{array}
    $$

\end{itemize}

\subsection*{2.3.5}
\begin{itemize}
    \item a - Valid \\
    \\
    $H(x):$ Student on the honor roll \\
    $A(x):$ Student got an A \\
    $D(x):$ Student got detention \\

    $$
    \begin{array}{c}
        \text{Form of the argument} \\
        \hline \\
        \forall x(H(x) \rightarrow A(x)) \\
        \forall x(D(x) \rightarrow \neg A(x)) \\
        \hline
        \therefore \forall x(D(x) \rightarrow \neg H(x))
    \end{array}
    $$
    \\

    \begin{center}
    Proof
    \end{center}
    
    $$
    \begin{array}{|c|c|c|}
        \hline
        1. & \forall x(H(x) \rightarrow A(x)) & \text{Hypothesis} \\
        \hline
        2. & c \text{ is an arbitray element in the domain} & \text{Element definition} \\
        \hline
        3. & H(c) \rightarrow A(c) & \text{Universal instantiation } 1,2 \\
        \hline
        4. & \forall x(D(x) \rightarrow \neg A(x)) & \text{Hypothesis} \\
        \hline
        5. & D(c) \rightarrow \neg A(c) & \text{Universal instantiation } 2,4 \\
        \hline
        6. & \neg H(c) \lor A(c) & \text{Conditional identity } 3 \\
        \hline
        7. & \neg D(c) \lor \neg A(c) & \text{Conditional identity } 5 \\
        \hline
        8. & \neg D(c) \lor \neg H(c) & \text{Resolution } 6,7 \\
        \hline
        9. & D(c) \rightarrow \neg H(c) & \text{Conditional identity } 8 \\
        \hline
        10. & \forall x(D(x) \rightarrow \neg H(x)) & \text{Universal generalization } 9 \\
        \hline
    \end{array}
    $$

    \item c - Invalid, $M(x)=F, D(x)=T$ \\
    \\
    $M(x):$ \text{The student missed class} \\
    $D(x):$ \text{The student got detention} \\

    $$
    \begin{array}{c}
        \text{Form of the argument} \\
        \hline \\
        \forall x(M(x) \rightarrow D(x)) \\
        Penelope \text{ is a particular element} \\
        D(p) \\
        \hline
        \therefore M(p)
    \end{array}
    $$

\end{itemize}

\subsection*{2.4.2}
\begin{itemize}
    \item b - Odd, $4n + 3 = 2(2n + 1) + 1$
    \item c - Even, $10n^3 + 8n - 4 = 2(5n^3 + 4n -2)$
\end{itemize}

\subsection*{2.4.6}
\begin{itemize}
    \item a - $x \ge 7$
    \item d - $x < 7$
\end{itemize}

\end{document}